%%
%% UTF-8       UTF-8      UTF-8       UTF-8       UTF-8       UTF-8
%% UTF-8       UTF-8      UTF-8       UTF-8       UTF-8       UTF-8
%%
%% XeLaTeX  XeLaTeX  XeLaTeX  XeLaTeX  XeLaTeX  XeLaTeX  XeLaTeX  XeLaTeX  

% !TeX program = xelatex
\lccode`\-=`\-
%\defaulthyphenchar=127
\documentclass[12pt,a4paper,oneside]{extarticle}
\usepackage{xunicode} % https://ctan.org/pkg/xunicode
\usepackage{xltxtra} % https://ctan.org/pkg/xltxtra
\usepackage{fontspec} % https://ctan.org/pkg/fontspec
\usepackage{polyglossia} % https://ctan.org/pkg/polyglossia
\setdefaultlanguage{portuges}
\usepackage{csquotes}

% Compilar com XeLaTeX
\usepackage{fontspec}
%\setmainfont{LMRoman10-Regular}
%\setsansfont{LMSans10-Regular}
%\setmonofont{LMMono10-Regular}

\clubpenalty=100
\linepenalty=100
\widowpenalty=100


\usepackage[backend=biber, style=numeric]{biblatex} 
\addbibresource{bibliografia.bib}

%estética
% \usepackage[hyphens]{url}
\usepackage{xurl}
\urlstyle{same} % URLs na mesma fonte do texto (Times New Roman)
\AtBeginDocument{%
  \setlength{\Urlmuskip}{0mu plus 1mu} % Ajusta espaçamento entre caracteres em URLs
  \setlength{\emergencystretch}{3em} % Permite maior flexibilidade na justificação
  \def\UrlBreakPenalty{100} % Penalidade baixa para facilitar quebras
  \def\UrlBigBreakPenalty{50} % Penalidade para quebras em caracteres como /
  \sloppy % Relaxa a justificação para evitar ultrapassagem
}

\usepackage[hang,small,bf]{caption}
\usepackage[a4paper,left=2.5cm,right=2.5cm,top=2.5cm,bottom=2.5cm,includehead,includefoot]{geometry}
\linespread{1.3}

\usepackage[
  xetex,                      % Driver para XeLaTeX
  unicode=true,               % Suporte a caracteres Unicode (ex.: ç, ã)
  bookmarks=true,             % Cria bookmarks no PDF
  bookmarksopen=true,         % Bookmarks abertos por padrão
  bookmarksnumbered=true,     % Numera os bookmarks (ex.: 1, 1.1, 1.2)
  linktoc=all,                % Links clicáveis em sumário, figuras e citações
  breaklinks=true,            % Permite quebra de links longos
  colorlinks=true,            % Links com cor em vez de caixas
  citecolor=black,            % Cor das citações (preto)
  filecolor=black,            % Cor dos links de arquivos (preto)
  linkcolor=black,            % Cor dos links internos (preto)
  urlcolor=blue,              % URLs em azul
  pdfpagelabels=true,         % Adiciona rótulos às páginas do PDF
  plainpages=false,           % Diferencia numeração de páginas
  pdftitle={Instalação do Arch Linux com XFCE no Android},
  pdfauthor={Luís A.V. Ferreira}, % Nome do autor
  pdfsubject={Instalação do Arch Linux com XFCE no Android},
  pdfkeywords={Arch Linux XFCE Android},
  pdfcreator={XeLaTeX com hyperref},
  pdfproducer={XeLaTeX}
]{hyperref}

%% Definir um novo estilo de url
\makeatletter
\def\url@luisstyle{%
  \@ifundefined{selectfont}{\def\UrlFont{\sf}}{\def\UrlFont{\small\ttfamily}}}
\makeatother
\urlstyle{luis}

\setlength{\captionmargin}{1.5in}
\usepackage{fancyhdr}
%\setlength{\headheight}{1in}

\renewcommand{\thefootnote}{\alph{footnote}} %modificar as letras de notas de rodapé%
\makeatletter
\renewcommand\@makefnmark{\@textsuperscript{\normalfont(\@thefnmark)}}
\makeatother
\renewcommand{\fboxrule}{2pt} %para alterar a expessura das caixas
%\renewcommand{\fboxsep}{3ex} %para alterar a separação nas caixas

\long\def\symbolfootnote[#1]#2{\begingroup\def\thefootnote{\fnsymbol{footnote}}\footnote[#1]{#2}\endgroup}

%%%%%%%%%%%%%%%%%%%%%%%%%%%%
%%%%%%%%%%%%%%%%%%%%%%%%%%%%
%%%%%%%%%%%%%%%%%%%%%%%%%%%%
\usepackage{setspace}
%\singlespacing
%\onehalfspacing
%\doublespacing
\usepackage{comment}
\usepackage{siunitx}
\usepackage{lettrine}
\usepackage{GoudyIn}
%%%%%%%%%%%%%%%%%%%%%%%%%%%%
\usepackage[x11names]{xcolor} 
\renewcommand{\LettrineFontHook}{\color{VioletRed4}\GoudyInfamily{}}
\LettrineTextFont{\itshape}
\setcounter{DefaultLines}{3}%
%%%%%%%%%%%%%%%%%%%%%%%%%%%%
\usepackage{relsize}
%\newcommand\CC{C\nolinebreak\hspace{-.05em}\raisebox{.4ex}{\relsize{-3}{\textbf{+}}}\nolinebreak\hspace{-.10em}\raisebox{.4ex}{\relsize{-3}{\textbf{+}}}}
\newcommand\CC{C\nolinebreak[4]\hspace{-.05em}\raisebox{.4ex}{\relsize{-3}{\textbf{++}}}}

%\usepackage{showframe}
\usepackage{pgf,tikz}       % - Elementos gráficos
\usepackage{tikzpagenodes}
\usetikzlibrary{backgrounds,calc,positioning}
\usetikzlibrary{positioning, shapes}

\pgfdeclarelayer{background}
\pgfdeclarelayer{foreground}
\pgfsetlayers{background,main,foreground}

\addbibresource{bibliografia.bib}

% \addbibresource{biblio.bib}
% \addbibresource{legislacao.bib}


\usepackage{listings}
\usepackage{lstlinebgrd}

% Definir cores personalizadas
\definecolor{lightgray}{gray}{0.95}
\definecolor{darkgray}{gray}{0.4}
\definecolor{purple}{rgb}{0.58,0,0.82}

% Traduzir o título "Listing" -> "Listagem"
\renewcommand{\lstlistingname}{Listagem}
\renewcommand{\lstlistlistingname}{Lista de Listagens}

% Definição do estilo
\lstset{
  basicstyle=\ttfamily\small,
  keywordstyle=\color{blue},
  stringstyle=\color{red!70!black},
  commentstyle=\color{green!50!black},
  breaklines=true,
  showstringspaces=false,
  upquote=true,
  extendedchars=true,
  literate={~}{{\textasciitilde}}1,
  numbers=left,                     % numeração das linhas à esquerda
  numberstyle=\tiny\color{darkgray},% estilo dos números
  numbersep=8pt,                    % distância entre número e código
  backgroundcolor=\color{lightgray},% cor de fundo base
  frame=single,                     % moldura à volta do código
  rulecolor=\color{black!30},       % cor da moldura
  language=python,                  % linguagem predefinida
  captionpos=t,                     % legenda em baixo
  xleftmargin=2em,                  % margens adicionais para espaço
  aboveskip=1em, belowskip=1em,
  linebackgroundcolor={%
        \ifodd\value{lstnumber}%
            \color{LightYellow1}%
        \else%
            \color{DarkSeaGreen1}%
        \fi}
}
% === Só isto resolve tudo (pythontex) ===
\usepackage{pythontex}
\setpythontexworkingdir{.}
\usepackage{float}      % <--- esta linha resolve o erro do [H]
 % não é para alterar

% Adicionar pacotes para controle de quebras
\usepackage{needspace}
\usepackage{afterpage}

\begin{document}

\begin{titlepage}
	\centering
	\includegraphics[width=0.15\textwidth]{logos/aeB06.pdf}  \par
	\vspace{1cm}
	{\scshape\LARGE Escola Secundaria José Gomes Ferreira    \par}
	\vspace{1cm}
	{\scshape\Large Física 12º     \par}
	\vspace{1.5cm}
	{\huge\bfseries Modelos de relatório}\par 
	\vspace{2cm}
	{\Large\itshape Adrian Alves Dias}\par
	{\Large\itshape Nº1  12º3}\par 
	\vfill
% Bottom of the page
	{\large \today\par}

	
\begin{tikzpicture}[remember picture, overlay]
\draw[line width = 2pt, color=black ] ($(current page.north west) + (.5in,-.5in)$) rectangle ($(current page.south east) + (-.5in,.5in)$);
\draw[line width = 1pt, color=ELECTRICRED ] ($(current page.north west) + (.55in,-.55in)$) rectangle ($(current page.south east) + (-.55in,.55in)$);
\draw[line width = 1pt, color=TANGERINYELLOW ] ($(current page.north west) + (.6in,-.6in)$) rectangle ($(current page.south east) + (-.6in,.6in)$);
\end{tikzpicture}

\end{titlepage}
 % pode alterar o título, o nome o número etc.

\begin{abstract}
Este trabalho apresenta a modelação e visualização 3D da topografia da região metropolitana do Rio de Janeiro utilizando \texttt{Python} e a biblioteca \texttt{pygmt}. Descreve-se a implementação computacional de um mapa de elevação com perspectiva tridimensional, abrangendo uma área de aproximadamente 30x25 km. O modelo desenvolvido incorpora dados de relevo terrestre de alta resolução, técnicas de sombreamento e representação de curvas de nível, proporcionando uma análise detalhada da configuração topográfica da região.
\end{abstract}

\section{Introdução}

\subsection{Enquadramento e Objectivos}
O presente trabalho tem como objectivo principal a integração entre Python (com a biblioteca PyGMT) e LaTeX (via PythonTeX) para geração automática de cartas topográficas. A selecção desta abordagem justifica-se pela sua aplicação em contextos académicos e científicos, onde a precisão cartográfica e a automatização de processos são factores determinantes.

A metodologia adoptada integra conceitos de processamento de dados geoespaciais, visualização cartográfica e composição tipográfica, visando não apenas a reprodução gráfica de cartas topográficas, mas também a compreensão dos princípios técnicos que governam a sua geração.

\subsection{Fundamentos Técnicos do PyGMT e PythonTeX}

\subsubsection{PyGMT e Processamento de Dados Geoespaciais}
O PyGMT constitui uma interface Python para o Generic Mapping Tools (GMT), oferecendo capacidades avançadas para manipulação de dados topográficos e geração de mapas de alta qualidade. A biblioteca suporta múltiplos formatos de dados, incluindo netCDF e GeoTIFF, e permite acesso a bases de dados públicas como o modelo global SRTM15+.

\subsubsection{Integração com LaTeX via PythonTeX}
O pacote PythonTeX permite a execução de código Python durante a compilação de documentos LaTeX, possibilitando a geração dinâmica de conteúdo. Esta abordagem garante reprodutibilidade e sincronização entre dados, cálculos e documentação.

\section{Objectivos do Trabalho}

\begin{itemize}
    \item Integrar código Python com PyGMT num documento LaTeX utilizando PythonTeX
    \item Executar cálculos e gerar automaticamente uma carta topográfica hipsométrica da região do Rio de Janeiro
    \item Utilizar dados SRTM com resolução de 1 segundo de arco (≈30 metros)
    \item Documentar todo o procedimento, incluindo código, figuras e referências bibliográficas
    \item Implementar visualização 3D com perspectiva realista utilizando \texttt{grdview}
\end{itemize}

\section{Metodologia}

\subsection{Abordagem Técnica}
A metodologia consistiu numa abordagem integrada que combina processamento de dados geoespaciais com composição documental automatizada. O fluxo de trabalho compreendeu as seguintes etapas principais:

\begin{itemize}
    \item Definição da região de estudo e parâmetros cartográficos
    \item Aquisição e processamento de dados topográficos via PyGMT
    \item Implementação de visualização 3D com configurações de perspectiva
    \item Integração do código Python no documento LaTeX via PythonTeX
    \item Processo de compilação em três fases para geração do documento final
\end{itemize}

\subsection{Configuração do Ambiente}
Para garantir reprodutibilidade, foram estabelecidas especificações técnicas precisas:

\begin{itemize}
    \item \textbf{PyGMT}: Versão 0.16.0 para compatibilidade com dados SRTM15+
    \item \textbf{Python}: Versão ≥3.10 para suporte às bibliotecas necessárias
    \item \textbf{LaTeX}: Distribuição TeXLive com XeLaTeX e pacote PythonTeX
    \item \textbf{Dados}: Modelo SRTM15+ com resolução de 1 segundo de arco
\end{itemize}

\section{Desenvolvimento}

\subsection{Definição da Área de Estudo}
A região seleccionada para análise corresponde à área metropolitana do Rio de Janeiro, delimitada pelas seguintes coordenadas geográficas:

\begin{itemize}
    \item \textbf{Longitude}: -43.35° a -43.05° (amplitude de 0.3°)
    \item \textbf{Latitude}: -23.05° a -22.83° (amplitude de 0.22°)
    \item \textbf{Área aproximada}: 30 × 25 km
    \item \textbf{Resolução espacial}: 1 segundo de arco (≈30 metros)
\end{itemize}

\subsection{Implementação da Visualização 3D}

\subsubsection{Configuração da Perspectiva}
A visualização tridimensional foi configurada com parâmetros optimizados para representação realista:

\begin{itemize}
    \item \textbf{Azimute}: 135° (vista do sudeste)
    \item \textbf{Elevação}: 30° (ângulo de visualização balanceado)
    \item \textbf{Projecção}: Mercator (M18c) com escala de 18 cm
    \item \textbf{Altura vertical}: 0 a 3500 metros
\end{itemize}

\subsubsection{Elementos Cartográficos}
Foram implementados múltiplos elementos para enriquecimento cartográfico:

\begin{itemize}
    \item \textbf{Coloração hipsométrica}: Paleta "geo" para representação natural do relevo
    \item \textbf{Curvas de nível}: Intervalo de 100m com anotações a cada 500m
    \item \textbf{Iluminação e sombreamento}: Configuração +a45+nt0.8 para realce do relevo
    \item \textbf{Elementos costeiros}: Linha de costa com resolução fina e água semi-transparente
\end{itemize}


\section{Resultados}

\subsection{Análise da Carta Topográfica Gerada}
A implementação produziu uma carta topográfica 3D de alta qualidade (600 DPI) que demonstra capacidades avançadas de visualização geoespacial. A figura resultante apresenta as seguintes características técnicas:

\begin{itemize}
    \item \textbf{Resolução espacial}: 1 segundo de arco (detalhamento de ≈30m)
    \item \textbf{Dimensões}: 18 cm de largura com proporções cartograficamente correctas
    \item \textbf{Elementos visuais}: Combinação de cores hipsométricas, curvas de nível e sombreamento
    \item \textbf{Anotações}: Escala de cores, título descritivo e elementos de orientação
\end{itemize}

\begin{figure}[h!]
\centering
\includegraphics[width=0.9\linewidth]{mapatopografico.png}
\caption{Carta topográfica 3D do Rio de Janeiro gerada automaticamente com PyGMT, mostrando o relevo da região com resolução de 30 metros.}
\label{fig:mapa_topografico}
\end{figure}

\subsection{Avaliação Técnica da Implementação}
A integração PyGMT + PythonTeX + LaTeX mostrou-se robusta para geração automatizada de documentação científica. O processo de compilação em três fases garantiu:

\begin{itemize}
    \item \textbf{Sincronização}: Actualização automática dos resultados com alterações no código
    \item \textbf{Qualidade gráfica}: Exportação vectorial com 600 DPI para impressão em escala A1
    \item \textbf{Reprodutibilidade}: Execução consistente em diferentes ambientes
\end{itemize}







\section{Conclusão}

\subsection{Síntese das Principais Conquistas}
O presente trabalho demonstra a viabilidade técnica da integração entre PyGMT e LaTeX via PythonTeX para produção automatizada de cartas topográficas. A abordagem desenvolvida permite:

\begin{itemize}
    \item Geração dinâmica de visualizações geoespaciais directamente em documentos académicos
    \item Manutenção da qualidade cartográfica profissional em formatos vectoriais
    \item Automatização completa do fluxo desde dados brutos até documentação final
\end{itemize}

\subsection{Limitações e Desafios Técnicos}
Foram identificados alguns desafios durante o desenvolvimento:

\begin{itemize}
    \item \textbf{Dependências complexas}: Requer configuração específica de múltiplas bibliotecas
    \item \textbf{Tempo de processamento}: Aquisição de dados SRTM pode demorar 30-60 segundos
    \item \textbf{Recursos computacionais}: Visualizações 3D exigem capacidade gráfica adequada
\end{itemize}

\subsection{Perspectivas Futuras}
Futuras evoluções do trabalho poderão incluir:

\begin{itemize}
    \item Extensão para outras projecções cartográficas e escalas
    \item Implementação de animações temporais para análise de processos geomorfológicos
    \item Integração com bases de dados adicionais (geologia, hidrografia, etc.)
    \item Desenvolvimento de templates reutilizáveis para diferentes regiões geográficas
\end{itemize}

\begin{lstlisting}[language=python]
import pygmt

region = [-43.35, -43.05, -23.05, -22.83]  # ~30x25 km
print("Baixando grid SRTM 1s (~30m) - pode demorar 30-60s na primeira vez...")
grid = pygmt.datasets.load_earth_relief(
    resolution="01s",
    region=region,
    registration="gridline"
)
fig = pygmt.Figure()
perspective = [135, 30]  # Vista clássica do Rio (sudeste, 30° de elevação)
fig.grdview(
    grid=grid,
    region=region + [0, 3500],
    projection="M18c",
    perspective=perspective,
    frame=["a", "z500f100"],
    zsize="4c",
    surftype="i",
    shading="+a45+nt0.8",
    cmap="geo",
    contourpen="0.3p,gray30"
)
fig.grdcontour(
    grid=grid,
    levels=100,                    # a cada 100m
    annotation="500",              # anotar a cada 500m
    pen="0.4p,gray10",
    perspective=perspective
)
fig.coast(
    resolution="f",
    shorelines=True,
    water="lightblue@80",          # água semi-transparente
    borders="1/0.6p,gray",
    perspective=perspective
)
fig.colorbar(
    perspective=perspective,
    frame=["a500f100", "x+lElevação (m)", "y+lm"],
    position="JBR+o1c/1c+w14c/0.8c+h"
)
fig.text(
    text="Rio de Janeiro - Topografia 3D (30 m)",
    position="TC",
    font="18p,Helvetica-Bold,white",
    offset="0/0.8c",
    perspective=perspective,
    fill="black@90"  # fundo semi-transparente
)

fig.savefig("rio_3D_perfeito.png", dpi=600)
print("Mapa 3D salvo como 'rio_3D_perfeito.png' (600 DPI - pronto para impressão A1!)")

fig.show()
\end{lstlisting}


\newpage

\begin{thebibliography}{9}


\end{thebibliography}

\label{fim}
\end{document}
